\documentclass[11pt,a4paper]{article}
\usepackage[utf8]{inputenc}
\usepackage[T1]{fontenc}
\usepackage[margin=0.75in]{geometry}
\usepackage{xeCJK} % 处理中文
\usepackage{enumitem}
\usepackage{titlesec}
\usepackage{hyperref}

% 格式设置
\titleformat{\section}{\large\bfseries}{}{0em}{}[\titlerule]
\titlespacing{\section}{0pt}{1.2ex}{0.8ex}
\setlist[itemize]{noitemsep, topsep=0pt}
\pagestyle{empty}

\begin{document}

% 个人信息
\begin{center}
    {\LARGE \textbf{柳宣泽}} \\
    \vspace{0.5em}
    手机：13175857611 | 邮箱：1463967532@qq.com | 籍贯：浙江台州
\end{center}

% 教育背景
\section{教育背景}
\textbf{南京理工大学} \hfill \textbf{自动化学院} \hfill \textbf{智能电网专业（国家特色）} \\
平均学分绩点：85.5 \hfill 绩点排名：1/40 \\
外语能力：CET-4（564分）、CET-6（504分） \\
主修课程：C语言程序设计 (93)、模拟电子线路 (90)、数字电子线路 (95)、控制工程基础（87）、智能电网信息技术（92）、高等数学1（87）、高等数学2（88）、人工智能基础（89） \\
学习获奖：国家奖学金；两次专业特等奖学金

% 实践经历
\section{实践经历}
\textbf{南京理工大学up70战队核心队员} \hfill 2023.12 -- 2024.06 \\
\textbf{南京理工大学up70战队 算法组组长} \hfill 2024.09 -- 现在 \\
负责技术研发，项目管理与培训

% 科研经历
\section{科研经历}

\textbf{MPPT最大功率点跟踪控制} \hfill 2025年9月-2026年4月 \\
\textit{指导老师：殷明慧 \hfill 第一负责人}
\begin{itemize}
    \item \textbf{项目简介：} 针对MPPT最优转矩法在面对湍流时跟踪失效问题进行改良
    \item \textbf{负责工作：} 
        \begin{enumerate}
            \item 通过搜集文献查找主流的方案有最优转矩，爬山法与叶尖速比法
            \item 在simulink仿真中搭建风电系统模型并验证了传统最优转矩法在面对大湍流时出现跟踪失效问题
            \item 分析跟踪失效原因，并提出 PI法 与 更改最优转矩曲线法 来解决跟踪失效问题
        \end{enumerate}
    \item \textbf{项目成果：} 完成实验平台效果验证，在大湍流下实现多30\%风能捕获率
\end{itemize}

\vspace{0.5em}

\textbf{MPC机器人底盘运动控制} \hfill 2026年3月-2026年4月 \\
\textit{第一负责人}
\begin{itemize}
    \item \textbf{项目简介：} 基于ROBOCON赛场下的小场景高速度的高性能舵轮底盘开发
    \item \textbf{负责工作：} 
        \begin{enumerate}
            \item 搭建 mujoco 舵轮仿真环境，搭建比赛地图，配置真实车体质量分布与更改关节响应时间
            \item 利用do-mpc框架，采用 绝对坐标 与 运动学 建模，完成定点控制。
            \item 增加 移动参考点 与 线性时变参数 来实现基于三次样条的路线跟踪控制，将状态空间改成 极坐标 来优化到达目标点附近的时候舵向容易突变的问题。
        \end{enumerate}
\end{itemize}

\vspace{0.5em}

\textbf{基于SLAM码盘融合的机器人高频定位系统研发} \hfill 2024年9月-2025年4月 \\
\textit{第一负责人}
\begin{itemize}
    \item \textbf{项目概述：} 针对ROBOCON赛场下，针对动态、复杂的定位需求开发SLAM轮式里程计点激光融合定位算法
    \item \textbf{负责工作：} 
        \begin{enumerate}
            \item 完成point\_lio voxel\_slam fast\_lio算法的复现，雷达驱动的开发与 docker 环境封装
            \item 通过 插值算法 实现正交里程计SLAM的 多传感器紧耦合定位, 采用 map->odom坐标 表示码盘偏移
            \item 基于点激光传感器完成SLAM 初始位置的纠正算法, 纠正后误差从10-20cm降到2cm内
        \end{enumerate}
    \item \textbf{项目成果：} robocon2025年全国二等奖
\end{itemize}

% 获奖情况
\section{获奖情况}
\begin{itemize}
    \item robocon颗粒归仓 \hfill 国家级一等奖 \hfill 2024年
    \item robocon飞身上篮 \hfill 国家级二等奖 \hfill 2025年
    \item 中国工程机器人大赛智慧物流 \hfill 国家级 亚军 \hfill 2025年
    \item 中国机器人技能大赛单项赛道 \hfill 国家级 冠军 \hfill 2024年
\end{itemize}

% 个人评价
\section{个人评价}
\begin{itemize}
    \item 熟悉STM32开发, vscode, keil, cmake工具链，能使用c/c++混合编程
    \item 熟悉ROS，git，docker，有大型项目开发经历
\end{itemize}

\end{document}